% ============================ 九、模型评价与推广 ============================
\section{模型的评价、改进与推广}

\subsection{模型的优点}

\textbf{（1）数据契约先行，三问互相印证。}本文没有把数据审计当作例行清洗，而是
把对数底数、量纲混用、检出限左删失、0--6 分制标定四件事一次性做实并写成可复算
报告。其价值不在于“干净”，而在于\textbf{三问共享同一套定义}：问题二的 OAV 与
问题三的缺陷约束用的是同一份浓度还原规则，问题三的感官画像与问题一的评价用的是
同一个预测器与同一套标定。若各问各自解释数据，结论会互相矛盾而无法自证。

\textbf{（2）把“不确定性”做成交付物的一部分。}问题一不止给出 50 维分数，还给
出基于重复间分裂一致性的预测区间，并实测其覆盖率。对使用者而言，“这一维的
3.2 分是可靠的，那一维的 2.8 分只能当参考”比一个精确到小数点后两位的数字有用
得多。

\textbf{（3）归因结论带“稳定性”与“化学合理性”双闸门。}问题二的 Top-3 不是
单次拟合的重要性排序，而是五路证据的秩融合，并用 bootstrap 入选频率与跨方法
Jaccard 双指标量化稳定性，再用 OAV、化学类别、代谢来源、模块一致性四条判据
论证化学合理性。这使得“关键化合物”的结论可被第三方复核，而不是只能相信作者的
模型。

\textbf{（4）迁移强度是构造性可解释的。}问题三把跨风格迁移写成一个带单调约束的
一维搜索，输出一个标量 $\lambda^\star$。它既有明确的物理含义（在缺陷不恶化的
前提下能迁移多远），又给出构造性保证（缺陷单调不增），不需要事后逐个样本检查。

\subsection{模型的不足}

\textbf{（1）感官标签的可解释上界。}感官分是品评小组均值，其本身含主观方差；
本文在 \S5.5 实测发现，标签区分度（标准差）与 $R^2$ 呈正相关，说明\textbf{部分
维度的低 $R^2$ 源于标签噪声而非模型不足}。本文没有、也无法把这一部分误差与
模型误差彻底分离——这是所有基于人工感官评分的研究的共同限制。

\textbf{（2）小样本风格的生成质量受限。}22 种风格中 Faro 仅 3 款、Brut 4 款。
本文用收缩协方差与风格嵌入让稀有风格向总体借力，但 \S7.6 的 Fréchet 距离显示
稀有风格的质量确实系统性差于主流风格。若要真正解决，需要扩充稀有风格的样本。

\textbf{（3）缺陷约束依赖感官代理模型。}反式-2-壬烯醛与陈酒花标志物在附件中
没有可用阈值，本文改用感官代理模型的预测值作为判据。这意味着缺陷约束的可靠性
部分继承自问题一模型的可靠性；若某款酒的缺陷维度预测偏差较大，约束就会失准。

\textbf{（4）共表达图只用了线性相关。}皮尔逊相关无法刻画“抑制型”或“阈值型”
交互。虽然问题二的五路证据中 SHAP 与互信息能部分补偿，但图结构本身仍是线性的。

\subsection{模型的改进方向}

\begin{enumerate}[label=(\arabic*),leftmargin=2.2em,itemsep=1pt]
\item \textbf{标签侧}：若同一款酒有多次重复品评记录，可用多层模型把“品评员
  随机效应”与“酒款真实差异”分离，从而给出 $R^2$ 的\textbf{信度上界}。
\item \textbf{结构侧}：把共表达图换成偏相关（精度矩阵）或非参数依赖（距离相关），
  可刻画条件独立关系，更接近“直接驱动”而非“共同上游”。
\item \textbf{生成侧}：用条件扩散模型替代归一化流，在样本量允许时可获得更好的
  多模态刻画；并引入“化学可行性投影”（非负浓度、元素守恒）作为硬约束。
\item \textbf{闭环侧}：本文的问题二给出候选化合物、问题三给出可生成谱，二者可
  合成一个\textbf{“目标感官 $\rightarrow$ 反向设计化合物谱”}的优化回路，这已是
  产品配方设计的范畴。
\end{enumerate}

\subsection{模型的推广}

本文的三层框架（\textbf{审计统一契约 $\rightarrow$ 多任务预测 $\rightarrow$ 图归因
$\rightarrow$ 约束生成}）不依赖啤酒这一具体对象，可直接迁移到其他“多组分
化学谱 $\rightarrow$ 感官/品质”的场景：葡萄酒与烈酒的风格指纹、茶叶与咖啡的
风味分型、香精香料的配方优化、乃至发酵过程的质量控制。其中问题二的“五路证据
秩融合 + 双稳定性指标 + 化学合理性四判据”尤其通用——凡是需要在“高维相关特征
+ 小样本”下给出可复核归因的场景，都可直接套用。
